\subsection{Уравнение теплового баланса}
%%%%%%%%%%%%%%%%%%%%%%%%%%%%%%%%%%%%%%%%%%%%%%%%%%%%
%%%%%%%%%%%%%%%%%%%%%%%%%%%%%%%%%%%%%%%%%%%%%%%%%%%%
\z{В калориметре смешали десять порций воды. Первая порция имела массу $m = 1\,\text{г}$ и температуру $t = 1\,\degree\text{C}$, вторая — массу $2m$ и температуру $2t$, третья — $3m$ и $3t$, и так далее, а десятая — массу $10m$ и температуру $10t$. Определите установившуюся температуру смеси. Потерями теплоты пренебречь.}
%
{$T_{\text{к}} = 7\,\degree\text{C}$}

%%%%%%%%%%%%%%%%%%%%%%%%%%%%%%%%%%%%%%%%%%%%%%%%%%%%
\z{Брусок, нагретый до $90\,\degree\text{C}$, опустили в калориметр с водой. При этом температура воды повысилась с $20\,\degree\text{C}$ до $40\,\degree\text{C}$. Какой станет температура воды в калориметре, если, не вынимая первого бруска, в неё опустить ещё один такой же брусок, нагретый до $70\,\degree\text{C}$? Потерями пренебречь.\newpage}
%
{$T = 47\,\degree\text{C}$}

%%%%%%%%%%%%%%%%%%%%%%%%%%%%%%%%%%%%%%%%%%%%%%%%%%%%
\z{В цилиндрический стакан налита вода до уровня $h_0 = 10\,\text{см}$ при температуре $t_0 = 0\,\degree\text{C}$. В стакан бросают алюминиевый шарик, вынутый из другого сосуда с водой, кипящей при температуре $t_{\text{к}} = 100\,\degree\text{C}$. При этом уровень воды повышается на $\Delta x = 1\,\text{см}$. Какой будет установившаяся температура в стакане? Удельные теплоёмкости воды и алюминия $c_{\text{в}}=4200\,\text{Дж/(кг$\cdot$\degree\text{C})}$ и $c_{\text{а}}= 920\,\text{Дж/(кг$\cdot$\degree\text{C})}$, плотности воды и алюминия $\rho_{\text{в}}=1000\,\text{кг/м}^3$ и $\rho_{\text{а}} = 2700\,\text{кг/м}^3$.}
%
{$T = 5.6\,\degree\text{C}$}

%%%%%%%%%%%%%%%%%%%%%%%%%%%%%%%%%%%%%%%%%%%%%%%%%%%%
\z{Имеется три стакана, содержащие массы $m$, $2m$ и $3m$ воды. В первом стакане вода холодная, во втором — горячая, в третьем — имеет некоторую промежуточную температуру. Из первого стакана берут ложку воды и переливают во второй, при этом температура воды во втором стакане уменьшается на величину $\Delta t$. Затем ложку воды из второго стакана переливают в третий, температура воды в третьем стакане возрастает на $\Delta t/2$. Затем ложку воды из третьего стакана переливают в первый. На сколько изменится при этом температура воды в первом стакане? Потерями тепла пренебречь.}
%
{$\text{На } \frac{\Delta T}{2}$}

%%%%%%%%%%%%%%%%%%%%%%%%%%%%%%%%%%%%%%%%%%%%%%%%%%%%
\z{В трёх одинаковых теплоизолированных сосудах находится одинаковое количество масла при комнатной температуре. Нагретый металлический цилиндр опустили в первый сосуд. После того как между цилиндром и маслом установилось тепловое равновесие, цилиндр перенесли во второй сосуд. После того как и там установилось равновесие, цилиндр перенесли в третий сосуд. На сколько градусов повысилась температура масла в третьем сосуде, если во втором она возросла на $5\,\degree\text{C}$, а в первом — на $20\,\degree\text{C}$?}
%
{$\Delta t_3 = \frac{(\Delta t_2)^2}{\Delta t_1} = 1.25\,\degree\text{C}$}

%%%%%%%%%%%%%%%%%%%%%%%%%%%%%%%%%%%%%%%%%%%%%%%%%%%%
\z{Для отопления комнаты по теплоизолированной трубе с площадью поперечного сечения $S_1 = 10\,\text{см}^2$ подавалась горячая вода со скоростью $v_1 = 0.48\,\text{м/с}$. При этом её температура на входе в батарею была равна $t_1 = 80\,\degree\text{C}$, а на выходе — $t_2 = 78\,\degree\text{C}$. Во время ремонта старую трубу заменили на новую с площадью поперечного сечения $S_2 = 8\,\text{см}^2$. Определите мощность батареи до замены трубы. \\
С какой скоростью $v_2$ должна двигаться по новой трубе вода, имеющая температуру $t_3=82\,\degree\text{C}$ на входе в батарею, чтобы температура воздуха $t_0=22\,\degree\text{C}$ в комнате осталась прежней? Плотность воды $\rho=1000\,\text{кг/м}^3$, удельная теплоёмкость воды $c=4200\,\text{Дж/(кг$\cdot$\degree\text{C})}$.}
%
{$P = v_1 S_1 \rho c (t_1 - t_2) = 4\,\text{кВт}; \quad v_2 = v_1 \frac{S_1 (t_1 - t_2)}{S_2 (t_3 - t_4)} = 0.2\,\text{м/с}, \quad t_4 = 76\,\degree\text{C}$}

%%%%%%%%%%%%%%%%%%%%%%%%%%%%%%%%%%%%%%%%%%%%%%%%%%%%
\z{Теплоёмкость некоторых материалов может зависеть от температуры. Рассмотрим брусок массы $m_1=1\,\text{кг}$, изготовленный из материала, удельная теплоёмкость которого зависит от температуры $t$ по закону $c(t)=c_1 (1 + \alpha t)$, где $c_1=1.4\cdot10^3\,\text{Дж/(кг$\cdot$\degree\text{C})}$ и $\alpha=0.014\,\degree\text{C}^{-1}$. Такой брусок, нагретый до температуры $t_1=100\,\degree\text{C}$, опускают в калориметр, в котором находится некоторая масса $m_2$ воды при температуре $t_2=20\,\degree\text{C}$. После установления теплового равновесия температура в калориметре оказалась равной $t_0= 60\,\degree\text{C}$. Пренебрегая теплоёмкостью калориметра и тепловыми потерями, определите массу $m_2$ воды в калориметре. Известно, что удельная теплоёмкость воды $c_2=4.2\cdot10^3\,\text{Дж/(кг$\cdot$\degree\text{C})}$.}
%
{$m_2 = 707\,\text{г}$}

 